\subsection*{The needy servant prays within Christ's body\enspace\properrefs{Int.}}

Psalm~85's opening appeal joins hearing, salvation, mercy and joy. The
antiphon selects clauses from verses~1--3; verse~4 then gives the servant's
lifting of the soul. The complete psalm names poverty and holiness in the
omitted clauses, widens into the nations' worship in verses~8--10, and
returns to threatening enemies in verses~14--17. Its account of God's
compassion, patience, mercy and truth in verse~15 recalls Exodus~34:6,
where divine forgiveness accompanies covenant renewal. The servant's
persistence rests on that character of God.

Augustine hears Christ together with his members praying through this
psalm (\emph{Enarrationes} 85.1--6, Tweed's English). Divine condescension
answers humility; the holiness for which the servant asks preservation is
itself God's grace. The Church's prayer extends across generations, and
love lifts the soul toward God. Some of this argument concerns the
poor-and-holy clauses beyond the antiphon's selection. Bellarmine expressly
receives Augustine's account of holiness through Christ and baptism and
explains the petition as a request for persevering salvation
(O'Sullivan's English, pp.~259--260).

Theodoret begins from David's prayer and Hezekiah's trust. Human nature is
poor even in a righteous person; the holiness claimed against unjust
attackers is relative innocence. His image of the inclined ear is a hearer
bending toward someone weakened by illness (PG~80:1553C--1556B).
Augustine's received holiness and Theodoret's relative innocence answer
different aspects of the complete psalm. Both direct weakness toward mercy.
At Mass the continuing cry stands beside the Collect's continuing mercy and
the Gospel's compassion for a mother whose request Luke never narrates.

\subsection*{Continuing mercy cleanses and fortifies the Church\enspace\properrefs{Coll.}}

The Collect makes continuing mercy the agent of cleansing and fortification.
Its reason is the Church's inability to stand safely without God; its
conclusion asks perpetual governance through God's gift. Stability thus
requires continuing divine action. The Epistle gives that preserved life
its communal tasks, and the Postcommunion asks the gift to work within
bodies and minds.

The prayer already belongs to a group with \latin{Tua nos} and
\latin{Mentes nostras} in Wilson's Gelasian III.xi, p.~230, and Gregorian
supplement XXXIIII, p.~174. Their changing headings and chant associations
receive fuller treatment below. Schuster's account of the Collect likewise
joins the Church's dependence to divine governance (III, p.~139).
The prayer asks both purification from what corrupts and strengthening for
what lies ahead; that second petition prepares the Secret's explicit plea
for protection.

\subsection*{The law of Christ makes correction a work of charity\enspace\properrefs{Ep.}}

Galatians~5:25--6:10 brings the letter's grace, faith and Spirit into
community conduct. The preceding discussion has made love and service the
form of freedom (5:6,13--24); the following conclusion returns to the cross
and new creation (6:11--18). Within the appointed passage, vainglory,
provocation and envy yield to gentle restoration, self-examination, support
of others, teaching, sowing and persevering good. The spiritual person must
consider his own susceptibility to temptation while correcting another.

Chrysostom calls the corrector's mildness itself a gift of the Spirit.
Reciprocal endurance fulfills Christ's law without requiring every gift or
function in every person; support of teachers cultivates humility on both
sides (\emph{Galatians}, chs.~5--6, lemmata 5:25--6:10,
Alexander's English). His explanation follows the whole appointed extent.
Theodoret's exposition of 6:1--10 gives restoration the form of care by
persons who know their own mutability. He identifies Christ's law with
charity and John~13:34, and at 6:10 invokes the heavenly Father's merciful
beneficence (PG~82:499A--502B). Chrysostom's allocation of different gifts
and Theodoret's shared mutability give reciprocal help complementary
reasons.

The commands to bear another's burden and one's own account belong to the
same argument. Augustine's \emph{Frangipane V} §§2--3 distinguishes shared
infirmity from personal reckoning, then counts neglect of a brother's wound
among one's own faults. The opening of §4 compares correction adapted to
weakness with a physician's treatment (Morin, pp.~213--215). Personal
answerability therefore intensifies the duty to help. Morin marks a repeated
charity formula at p.~214 as a possible marginal gloss; the surrounding
argument already joins charity and responsibility without that repetition.

Aquinas specifies three forms of burden-bearing: enduring bodily or
spiritual defects, relieving actual need, and offering satisfaction through
prayer and good works (\emph{Super Galatas} 6.1, Latin electronic unit
[87781]). These are distinct acts rather than a vague disposition of
sympathy. The personal account of 6:5 survives all three. In lecture~2
[87782], instruction receives material support, sowing takes its life from
the Spirit who comes from God, and universal beneficence rests upon the
common divine image. Particular care for fellow believers orders practical
obligation within good done to all.

Paul's harvest makes present opportunity serious. Sowing to disordered
flesh yields corruption, while sowing to the Spirit yields eternal life;
perseverance awaits its fitting season. The Gospel's immediate rising and
this future harvest describe different gifts of life. The Postcommunion
asks for divine efficacy within the very agents whom Paul commands to act.

\subsection*{Mercy and truth sustain praise through changing conditions\enspace\properrefs{Grad.}}

Psalm~91 opens with the goodness of praise, mercy in the morning and truth
by night. Its Sabbath heading and later contrast of passing wicked
prosperity with lasting righteous fruitfulness frame those lines
(vv.~1--16). The antiphon appoints the praise, while planting, cedar and
fruitful old age belong to the full psalm's development.

Augustine receives the Sabbath as the heart's rest and morning/night as
prosperity/adversity: mercy is praised amid good fortune, and just
correction trusted amid affliction (91.1--4, Tweed's English).
Theodoret instead makes night and day the comprehensive span of praise for
beneficence and righteous judgment, within a Sabbath that points toward
future rest (PG~80:1615C--1618A). Bellarmine first distinguishes mercy
manifest in creation from judgment praised in faith's darkness, then
expressly takes up Augustine's account (O'Sullivan, pp.~289--290).
The difference matters when prosperity and apparent injustice coexist:
Bellarmine addresses what can be seen, Augustine the conditions endured,
and Theodoret the praise offered continually. Their readings strengthen
the Epistle's perseverance without identifying an immediate favorable
outcome with the eternal harvest.

\subsection*{The great King rules all the earth\enspace\properrefs{All.}}

The Alleluia's \latin{rex magnus super omnem terram} has an exact verbal
parallel in Clementine Psalm~46:3. Psalm~94:3, the Missal's principal
reference, ends \latin{super omnes deos}; Challoner renders that ending
``above all gods.'' The full Psalm~94 passes from creation and shepherding
to Israel's hardened hearts in the wilderness, the passage Hebrews~3--4
receives in its teaching about hearing and rest. The chant selects the
royal acclamation.

Augustine's Psalm~94 lemma has the all-gods form. He places kingship within
Israel's faithful remnant and, at the following verse, Christ the
cornerstone's union of circumcision and uncircumcision (94.5--6,10).
Bellarmine also has all gods and develops universal dominion in the next
verse (pp.~298--299). Theodoret's actual Psalm~94 lemma has all the earth,
in both Greek and facing Latin, and opposes universal rule to falsely
named gods (PG~80:1641A--1642A). This is a received textual analogue to
the Missal's form.

Augustine's separate Psalm~46.3 exposition receives the parallel all-earth
clause as the kingship of Christ crucified and ascended. Believing Jews,
including the apostles, stand within that dominion together with the nations
(Tweed's English, 46.2--3). The appointed acclamation thus gives the
Gospel's particular mercy a universal horizon. The Psalm~46 reception and
Theodoret's Psalm~94 lemma are distinct witnesses; neither supplies a
history of the chant's adaptation.

Theodoret's separate Psalm~46 exposition makes the invitation to praise
apostolic: the highest King's sovereignty, known in Israel, is manifested
to all nations (PG~80:1205D--1208B). Bellarmine explains the same royal
clause through Christ's divine greatness, irresistible power and
providential government of the world, which give all peoples reason for
joy (O'Sullivan's English, Psalm~XLVI, edition v.~2, p.~144).
Aquinas distinguishes universal dominion, eternal duration and authority
over kings; God's transcendence does not exclude his care of earthly
affairs (\emph{Super Psalmo}~46, n.~2, Latin electronic unit [87252]).
These direct expositions of the parallel clause give universal kingship
concrete force: the Lord whose mercy meets one mother also governs the
whole earth. Bellarmine's translated exposition omits philological
discussion; Aquinas's Latin is a modern electronic transcription.

\subsection*{Christ returns a living son to his mother\enspace\properrefs{Gosp.}}

Luke places Naim between the centurion's requested healing and John's
inquiry, answered in part by the raising of the dead (7:1--23).
Christ's followers meet the city's mourners at the gate. Jesus sees a
widow accompanying her only son's body, is moved with compassion, touches
the bier and commands the youth to arise. The youth sits up and speaks;
Jesus gives him to his mother. The crowd fears and glorifies God,
acclaiming a great prophet and God's visitation (7:11--16). Luke supplies
no request, confession of faith, imputation of sin, or report of the son's
words in this passage.

Elias returns a living son to a widow after petitioning God
(3~Kings~17:17--24); Eliseus prays before restoring the Sunamitess's child
(4~Kings~4:8--37). The latter mother has a living husband and is distinct
from the indebted widow at that chapter's opening. Jesus' own command and
return of the son give the parallels their particular force at Naim.
The crowd's visitation language also resonates with the Benedictus and
Exodus~4:29--31, where divine attention to affliction elicits worship.
Luke calls Jesus Lord in his narration; the crowd supplies the two
acclamations.

Cyril's \emph{Luke} Sermon~XXXVI explains command and touch through the
flesh united to the life-giving Word, and receives the miracle as a pledge
of the general resurrection (Payne Smith, pp.~132--135). That English
witness has a missing Syriac folio supplemented from a Greek catena at
p.~132 and additional Greek fragments in p.~135's note. The surviving
argument joins Christ's flesh to the Word's life-giving power; its
iron-in-fire comparison describes participation in that power without
destroying the flesh's humanity. Cyril ends with deliverance from evil
works and communion with the saints.

Augustine's complete \emph{Sermon}~98.1--7 (MacMullen's NPNF
Sermon~48) holds the bodily raising together with restoration from sin.
The three Gospel raisings distinguish inward consent, public action and
hardened habit; Naim supplies the public-action case at §§4--5.
The widow is Mother Church receiving a restored child. Augustine's moral
schema interprets the miracle; Luke has charged this young man with no
sin.

Ambrose, in the continuous Latin transcription of his commentary, affirms
the literal compassion for the widow before opening the ecclesial reading. Her tears become
Mother Church's intercession; the touched wood signifies salvation through
the cross, and the Church intercedes for each sinner as for an only son
(\emph{Expositio in Lucam} V.89--92). Bede, in the continuous Latin
transcription of his commentary, develops conscience awakened, compunction and visible renewed
life. Return to the mother explicitly becomes reunion with the Church's
communion through priestly judgment (II.7, PL~92:418A--D). His argument
against denying penitents hope gives reconciliation an ecclesial form.
The return to communion links the Gospel's reception to Paul's
restorative charity and the Communion's abiding in Christ.

\subsection*{Patient petition receives a new canticle\enspace\properrefs{Off.}}

The Offertory selects waiting, divine regard, prayer heard and a canticle
placed in the mouth. Its substitutions of \latin{respéxit me},
\latin{deprecatiónem meam} and \latin{hymnum} differ from the Clementine
forms. The pit, mire, rock, steadied steps and many who see are omitted
parts of the complete Psalm~39. That psalm moves onward through obedience
and public testimony before renewing petition amid surrounding evils
(vv.~2--18).

Augustine hears Christ as Head with renewed members (39.2--5).
Rescue from sinful mire establishes the person on Christ the rock; the new
song belongs to renewed life, whose walk and future crown still lie ahead.
Theodoret reports other allocations of the pit before giving his own
Davidic and human-wide interpretation: the incarnate Savior rescues from
death and sin, and new benefits receive a new song
(PG~80:1151B--1154C). The pit imagery is the full-psalm context.
Bellarmine reads Christ's people awaiting redemption and established in
faith, hope and charity (pp.~118--119). Thanksgiving belongs to a people
who remain dependent upon the Savior.

Hebrews~10:5--10 receives the psalm's later obedience passage in relation
to Christ's bodily offering. Its body wording differs from the ears wording
in Clementine Psalm~39:7. That canonical development concerns an unappointed
portion of the psalm. At this Sunday's Offertory, \emph{The Liturgical
Year} explicitly joins Mother Church's restored children to thanksgiving
(1909, p.~354), drawing the Gospel's ecclesial reception into the assembly's
song.

\subsection*{The sacraments guard against assault\enspace\properrefs{Sec.}}

The Secret asks the Lord's sacraments to guard the faithful and defend them
continually against demonic assaults. The enemy named is demonic; the
prayer does not redefine human opponents as demons. The Epistle's own
vulnerable correctors need strength for the charity they undertake.
The 1962 Vatican witness prints the special conclusion through
\latin{in unitáte} and appoints the Trinity Preface; Benziger abbreviates
that conclusion, while the Vatican text controls this formulary.

Aquinas gives protection a twofold account: inward union with Christ
strengthens spiritual life, while the Passion defends against hostile
attack (\emph{Summa} III.79.6). A recipient can still sin. Protection
therefore sustains vigilance and conversion rather than removing their
necessity. Schuster joins custody to nourishment and the faculties governed
by grace (III, pp.~139--141). The Secret asks for defense within that
living relation to Christ.

\subsection*{Christ gives his flesh for life\enspace\properrefs{Comm.}}

The Communion appoints the last clause of Vulgate John~6:52, corresponding
to modern 6:51b. It says \latin{dédero} and \latin{sæculi}, where
Clementine John has \latin{dabo} and \latin{mundi}. The complete chapter
moves from feeding and travel to manna, heavenly teaching, Christ's flesh,
abiding, the refusal of many disciples and Peter's confession.
The discourse explicitly contrasts the ancestors who ate manna and died
with life and resurrection through Christ (6:49--60).

Cyril explains the gift through the flesh united to the life-giving Word:
Christ offers his flesh for the world, and his death defeats death
(\emph{John} IV.2, Pusey's English, pp.~409--414).
At pp.~418--419, on a following verse, Cyril expressly invokes the
life-giving touch at Naim to explain the greater power of sacramental
participation. His distinction between the resurrection of all and the
blessed life of the holy prevents world-directed gift from becoming an
assertion of automatic beatitude.

Augustine's Tractate~26.10--20 explains unity and charity within Christ's
body. Section~13 joins the body animated by Christ's Spirit to one bread
and one body; §§15 and 18 distinguish the visibly received sacrament from
holy fellowship and abiding in Christ. Section~16 retains future bodily
resurrection despite temporal death. Outward eating alone does not secure
that saving reality (Gibb's English).

Chrysostom's Homily~46.1 permits doctrine/faith or Christ's body as readings
of earlier bread-of-life language. At §§2--3 he develops the flesh,
manna's fulfillment and real incorporation through the offered food;
§4 demands restraint of wrath and speech and works accompanying faith
(Marriott's English). The earlier alternatives and the later sacramental
exposition have different verse loci. Aquinas likewise distinguishes
figure from contained reality and sacramental sign from spiritual eating
(\emph{Super Ioannem} 6.6, Latin unit [87463]). The Passion applies to
present spiritual life and future fulfillment within ecclesial unity.
The Communion names the gift whose effects the Postcommunion then asks to
possess the whole person.

\subsection*{The heavenly gift possesses minds and bodies\enspace\properrefs{Postcomm.}}

The final prayer gives the operation of the heavenly gift a double object:
minds and bodies. Its effect is to precede \latin{noster sensus} within
the recipients. The Church who could not stand safely unaided now asks
for divine action within her members, sustaining the conduct Paul has
required. Schuster interprets the prayer as the faculties' inclinations
brought under grace (III, p.~141); \emph{The Liturgical Year} joins this
mind-and-body efficacy to the life-giving Communion (1909, p.~355).

Aquinas distinguishes bodily conduct receiving the gift's effect now from
incorruption and glory in the future (III.79.1, ad~3). Charity is stirred
to act (ad~2), while spiritual nourishment presupposes the living relation
which knowingly retained mortal sin obstructs (III.79.3).
The prayer's comprehensive reach embraces embodied action today and its
promised fulfillment. Christ's life-giving flesh sustains the Church's
patient, responsible life rather than leaving the body outside salvation.
